% Part: lambda-calculus
% Chapter: syntax
% Section: alpha

\documentclass[../../../include/open-logic-section]{subfiles}

\begin{document}

\olfileid[fr]{lam}{syn}{alp}
\olsection{\texorpdfstring{$\alpha$}{α}-conversion}

Quel rapport y a-t-il entre $\lambd[x][x]$
et $\lambd[y][y]$? Les deux termes représentent
la fonction identité, bien qu'ils soient
syntaxiquement distincts. Ils ne diffèrent
que par le nom de leur variable liée : on
obtient l'un en « renommant » la variable liée
de l'autre. C'est ce qu'on appelle une
\emph{$\alpha$-conversion}.

\begin{defn}[Changement de variable liée, $\aconvone$]
Si un terme $M$ contient une occurrence de
$\lambd[x][N]$, si $x\neq y$,
$y \notin \FV{N}$ et si $\Subst{N}{y}{x}$
est définie, remplacer cette occurrence par
\begin{equation*}
  \lambd[y][\Subst{N}{y}{x}]
\end{equation*}
donne un terme $M'$ et constitue un
\emph{changement de variable liée}, noté
$M \redone[\alpha] M'$.
\end{defn}

\begin{explain}
Le changement élémentaire exige deux noms distincts, même si
le corps ne contient aucune occurrence libre de la variable
du lieur. La conversion en zéro changement donnera ensuite
la réflexivité de l'$\alpha$-conversion.
\end{explain}

\begin{defn}[Compatibilité d'une relation]
Une relation $R$ sur les termes est dite
\emph{compatible} si elle satisfait les
conditions suivantes :
\begin{enumerate}
\item Si $R N N'$, alors
  $R \lambd[x][N] \lambd[x][N']$.
\item Si $R P P'$, alors $R (PQ) (P'Q)$.
\item Si $R Q Q'$, alors $R (PQ) (PQ')$.
\end{enumerate}
\end{defn}

Reformulons donc la définition.
\begin{defn}[Changement de variable liée, $\aconvone$]
Le \emph{changement de variable liée}
($\redone[\alpha]$) est la plus petite
relation compatible sur les termes qui
satisfait aussi la condition suivante :
\begin{align*}
  &\lambd[x][N] \redone[\alpha]
   \lambd[y][\Subst{N}{y}{x}]
   && \text{si $x \neq y$ et $y \notin \FV{N}$}\\
  & && \text{et si $\Subst{N}{y}{x}$ est définie.}
\end{align*}
\end{defn}

« La plus petite » signifie ici qu'elle ne
contient que les couples imposés par la
compatibilité et par cette condition. On
peut donc aussi la définir par induction.

\begin{defn}[Changement de variable liée, $\aconvone$]
\ollabel{defn:aconvone}
Le \emph{changement de variable liée}
($\aconvone$) est défini inductivement
comme suit :
\begin{enumerate}
\item Si $N \aconvone N'$, alors
  $\lambd[x][N] \aconvone \lambd[x][N']$.
  \ollabel{defn:aconvone1}
\item Si $P \aconvone P'$, alors
  $(PQ) \aconvone (P'Q)$.
  \ollabel{defn:aconvone2}
\item Si $Q \aconvone Q'$, alors
  $(PQ) \aconvone (PQ')$.
  \ollabel{defn:aconvone3}
\item Si $x \neq y$, $y \notin \FV{N}$ et
  $\Subst{N}{y}{x}$ est définie, alors
  $\lambd[x][N] \redone[\alpha]
  \lambd[y][\Subst{N}{y}{x}]$.
  \ollabel{defn:aconvone4}
\end{enumerate}
\end{defn}

Ces définitions sont équivalentes; nous
laissons leur démonstration en exercice.
Dans la suite, nous utiliserons la
définition inductive.

\begin{defn}[$\alpha$-conversion, $\aconv$]
L'\emph{$\alpha$-conversion} ($\aconv$)
est la plus petite relation réflexive et
transitive sur les termes qui contient
$\aconvone$.
\end{defn}

Comme ci-dessus, « la plus petite » signifie
qu'on n'ajoute que les couples imposés par
la réflexivité, la transitivité et
$\aconvone$. On obtient ainsi la définition
équivalente suivante.

\begin{defn}[$\alpha$-conversion, $\aconv$]
\ollabel{defn:aconv}
L'\emph{$\alpha$-conversion} ($\aconv$)
est définie inductivement comme suit :
\begin{enumerate}
\item Si $P \aconv Q$ et $Q \aconv R$,
  alors $P \aconv R$.
  \ollabel{defn:aconv1}
\item Si $P \aconvone Q$, alors
  $P \aconv Q$.
  \ollabel{defn:aconv2}
\item $P \aconv P$.
  \ollabel{defn:aconv3}
\end{enumerate}
\end{defn}

\begin{ex}
$\lambd[x][f x]$ se $\alpha$-convertit en
$\lambd[y][f y]$, et réciproquement. De
manière informelle, ces deux termes désignent
une fonction qui reçoit un argument et renvoie
le résultat de l'application de~$f$ à celui-ci;
$f$ est une variable de l'environnement.

$\lambd[x][f x]$ ne se $\alpha$-convertit pas
en $\lambd[x][g x]$. Ces termes renvoient
respectivement aux variables d'environnement
$f$ et~$g$ : ils peuvent donc se comporter
différemment dans un environnement où
$f$ et $g$ ont des valeurs différentes.
\end{ex}

\begin{prob}
Les couples de termes suivants sont-ils
$\alpha$-convertibles?
\begin{enumerate}
\item $\lambd[x][\lambd[y][x]]$ et
  $\lambd[y][\lambd[x][y]]$
\item $\lambd[x][\lambd[y][x]]$ et
  $\lambd[c][\lambd[b][a]]$
\item $\lambd[x][\lambd[y][x]]$ et
  $\lambd[c][\lambd[b][a]]$
\end{enumerate}
\end{prob}

\begin{explain}
Les deux derniers items de l'exercice portent sur le même
couple de termes.
\end{explain}

\begin{lem}\ollabel{lem:fv-one}
Si $P \aconvone Q$, alors
$\FV{P} = \FV{Q}$.
\end{lem}

\begin{proof}
Raisonnons par induction sur la \emph{dérivation}
de $P \aconvone Q$.
\begin{enumerate}
\item Si la dernière règle est
  \olref{defn:aconvone4}, alors $P$ est de
  la forme $\lambd[x][N]$ et $Q$ de la forme
  $\lambd[y][\Subst{N}{y}{x}]$, où
  $x \neq y$, $y \notin \FV{N}$ et
  $\Subst{N}{y}{x}$ est définie.
  Distinguons deux cas selon que
  $x \in \FV{N}$ ou non :
  \begin{enumerate}
  \item Si $x \in \FV{N}$, alors :
    \begin{align*}
      \FV{\lambd[y][\Subst{N}{y}{x}]}
      &= \FV{\Subst{N}{y}{x}} \setminus \{y\}\\
      &= ((\FV{N} \setminus \{x\}) \cup \{y\})
         \setminus \{y\}
         && \text{par le \olref[sub]{thm:infv}}\\
      &= \FV{N} \setminus \{x\}\\
      &= \FV{\lambd[x][N]}.
    \end{align*}
  \item Si $x \notin \FV{N}$, alors :
    \begin{align*}
      \FV{\lambd[y][\Subst{N}{y}{x}]}
      &= \FV{\Subst{N}{y}{x}} \setminus \{y\}\\
      &= \FV{N} \setminus \{y\}
         && \text{par le \olref[sub]{thm:notinfv}}\\
      &= \FV{N} \setminus \{x\}
         && \text{car $x,y\notin\FV{N}$}\\
      &= \FV{\lambd[x][N]}.
    \end{align*}
  \end{enumerate}
\item Les trois autres cas sont laissés
  en exercice.
\end{enumerate}
\end{proof}

\begin{explain}
Si $x$ et $y$ sont absentes des variables libres du corps,
retirer l'une ou l'autre de cet ensemble le laisse inchangé.
C'est le cas du changement de nom d'un lieur sans occurrence
à lier dans son corps.
\end{explain}

\begin{prob}
Compléter la preuve du
\olref[lam][syn][alp]{lem:fv-one}.
\end{prob}

\begin{lem}\ollabel{lem:inv}
Si $P \aconvone Q$, alors
$Q \aconvone P$.
\end{lem}

\begin{proof}
Raisonnons par induction sur la \emph{dérivation}
de $P \aconvone Q$.
\begin{enumerate}
\item Si la dernière règle est
  \olref{defn:aconvone4}, alors $P$ est de
  la forme $\lambd[x][N]$ et $Q$ de la forme
  $\lambd[y][\Subst{N}{y}{x}]$, où
  $x \neq y$, $y \notin \FV{N}$ et
  $\Subst{N}{y}{x}$ est définie.
  D'après le \olref[sub]{thm:clr}, on a
  $x \notin \FV{\Subst{N}{y}{x}}$.
  D'après le \olref[sub]{thm:inv}, la substitution
  $\Subst{\Subst{N}{y}{x}}{x}{y}$ est définie
  et vaut~$N$. Dès lors,
  \olref{defn:aconvone4} donne
  $\lambd[y][\Subst{N}{y}{x}]
  \aconvone \lambd[x][\Subst{\Subst{N}{y}{x}}{x}{y}]
  = \lambd[x][N]$.
\end{enumerate}
\end{proof}
\begin{explain}
Pour revenir au nom~$x$, il faut que $x$ soit absent des
variables libres du corps obtenu. Par exemple, si $N=x$,
le nouveau corps est~$y$ et contient $y$ libre.
\end{explain}
\begin{prob}
Compléter la preuve du
\olref[lam][syn][alp]{lem:inv}.
\end{prob}

\begin{thm}
L'$\alpha$-conversion est une relation
d'équivalence sur les termes : elle est
réflexive, symétrique et transitive.
\end{thm}

\begin{proof}
\begin{enumerate}
\item Chaque terme $M$ peut être transformé
  en lui-même par \emph{zéro} changement
  de variable liée.
\item Si $P$ se $\alpha$-convertit en $Q$
  par une suite de changements de variables
  liées, on peut, à partir de $Q$, inverser
  ces changements grâce au \olref{lem:inv},
  dans l'ordre opposé, pour retrouver~$P$.
\item Si $P$ se $\alpha$-convertit en $Q$
  par une suite de changements, puis $Q$
  en $R$ par une autre suite, leur
  concaténation convertit $P$ en $R$.
\end{enumerate}
\end{proof}

Désormais, nous dirons que $M$ et $N$ sont
\emph{$\alpha$-équivalents}, $M \aeq N$,
si $M$ se $\alpha$-convertit en $N$;
par le théorème précédent, cela équivaut
à l'$\alpha$-conversion de $N$ en~$M$.

\begin{thm}\ollabel{thm:fv}
Si $M \aeq N$, alors $\FV{M} = \FV{N}$.
\end{thm}

\begin{proof}
Cela découle immédiatement du
\olref{lem:fv-one}.
\end{proof}

\begin{lem}\ollabel{lem:sub:R}
Si $R \aeq R'$ et si $\Subst{M}{R}{y}$
est définie, alors $\Subst{M}{R'}{y}$
est définie et $\alpha$-équivalente à
$\Subst{M}{R}{y}$.
\end{lem}

\begin{proof}
Exercice.
\end{proof}

\begin{prob}
Démontrer le \olref[lam][syn][alp]{lem:sub:R}.
\end{prob}

Rappelons que, dans la \olref[sub]{sec}, la
substitution est indéfinie dans certains cas.
La conversion des termes à une forme
$\alpha$-équivalente permet cependant de
renommer des variables liées afin de rendre
la substitution définie. Le théorème suivant
affirme aussi l'indépendance du résultat,
à $\alpha$-équivalence près.

\begin{thm}\ollabel{thm:sub}
Pour tous $M$, $R$ et~$y$, il existe un terme
$M'$ tel que $M \aeq M'$ et que
$\Subst{M'}{R}{y}$ soit définie. De plus,
si $M'' \aeq M$ et si $R''$ vérifie
$R'' \aeq R$ avec $\Subst{M''}{R''}{y}$
définie, alors
$\Subst{M'}{R}{y} \aeq
\Subst{M''}{R''}{y}$.
\end{thm}

\begin{proof}
Commençons par une observation sur l'$\alpha$-équivalence.
Dans l'arbre d'un terme, gardons les nœuds d'application et
d'abstraction et effaçons les noms portés par les abstractions.
À chaque occurrence de variable, inscrivons son nom si elle est
libre. Si elle est liée, inscrivons le nombre d'abstractions à
franchir jusqu'à son lieur, en comptant à partir de zéro pour
l'abstraction la plus proche.

Un changement de variable liée conserve cet arbre. En effet, dans
la \olref{defn:aconvone}, la clause \olref{defn:aconvone4}
change exactement les occurrences liées par l'abstraction visée.
Les conditions sur le nouveau nom et la substitution interdisent
toute capture. La clause \olref{defn:aconvone1} permet
d'effectuer ces changements sous une autre abstraction. Les deux
clauses d'application permettent de les effectuer dans chaque
branche. Il s'ensuit que des termes $\alpha$-équivalents ont le
même arbre ainsi étiqueté.

Réciproquement, supposons que deux termes aient le même arbre.
Associons à chaque position d'abstraction un nom distinct, absent
des deux termes. Dans chacun d'eux, parcourons les abstractions de l'intérieur
vers l'extérieur et renommons chaque lieur par le nom choisi
pour sa position. Au moment de renommer un lieur, tous ceux de
son corps portent déjà des noms frais et distincts. Aucun ne
porte donc l'ancien nom du lieur visé. Le nouveau nom est encore
absent du terme. La substitution de renommage est ainsi définie,
par induction sur le corps de l'abstraction. Après ce parcours,
les deux termes sont identiques.
En effet, les occurrences libres ont gardé leur nom et chaque
occurrence liée porte le nouveau nom du lieur indiqué par son
étiquette. Les deux termes sont donc $\alpha$-équivalents.

Pour l'existence, appliquons le même procédé à $M$, en choisissant
les noms des lieurs en dehors des noms de $M$, de $\FV{R}$ et de
$\{y\}$. Le résultat $M'$ est $\alpha$-équivalent à $M$. Tout
lieur de $M'$ est distinct de $y$ et absent de $\FV{R}$. Les
clauses de la substitution donnent donc, par induction sur $M'$,
que $\Subst{M'}{R}{y}$ est définie.

Pour l'unicité, considérons une substitution définie
$\Subst{U}{V}{y}$. Son arbre étiqueté s'obtient en remplaçant
chaque feuille libre de nom $y$ de l'arbre de $U$ par l'arbre
étiqueté de $V$. Cela se vérifie par induction sur $U$. À une
abstraction, la condition sur le lieur assure que les variables
libres du terme inséré restent libres. Les distances des
occurrences liées dans ce terme restent les mêmes, puisque
leurs lieurs se trouvent à l'intérieur du terme inséré. Comme
$M' \aeq M''$ et $R \aeq R''$, les deux substitutions de
l'énoncé donnent le même arbre étiqueté. La réciproque établie
plus haut donne leur $\alpha$-équivalence.
\end{proof}

\begin{prob}
Dans la preuve du \olref[lam][syn][alp]{thm:sub}, détailler
l'induction qui décrit l'arbre du résultat d'une substitution
définie. Traiter le cas d'une variable, celui d'une application
et celui d'une abstraction.
\end{prob}

\begin{cor}\ollabel{cor:sub}
Pour tous $M$, $R$ et~$y$, il existe un couple
$M'$, $R'$ tel que $M' \aeq M$,
$R \aeq R'$ et que $\Subst{M'}{R'}{y}$
soit définie. De plus, si un autre couple
$M'' \aeq M$, $R'' \aeq R$ vérifie
$\Subst{M''}{R''}{y}$ définie, alors
$\Subst{M'}{R'}{y} \aeq
\Subst{M''}{R''}{y}$.
\end{cor}

\begin{proof}
Pour l'existence, prenons $R'=R$ et appliquons le
\olref{thm:sub}. Pour l'unicité, appliquons ce même
théorème au terme de départ $M$ et au couple $M'',R''$.
\end{proof}

\end{document}
